\documentclass[12pt]{article}
\usepackage{amsmath,amssymb,amsfonts,amsthm}
\usepackage[margin=1in]{geometry}

\begin{document}

\begin{center}
{\Large\bfseries Chapter 9, Problem 6}\\[6pt]
Byron \& Fuller, \textit{Mathematics of Classical and Quantum Physics}
\end{center}

\bigskip

\textbf{Problem.} Let $L_2[0,1]$ denote the space of square-integrable functions on $[0,1]$. Define the operator $P$ on $L_2[0,1]$ by $Pf(t) = f(1-t)$, and the integral operator $A$ with kernel $a(t,t') = t_<(1-t_>)$, where $t_<$ and $t_>$ denote the smaller and larger of $t$ and $t'$.

(a) Show that $P$ commutes with $A$.
(b) Show that $P$ is self-adjoint and bounded. Find $\|P\|$.
(c) Find the possible eigenvalues of $P$. Is $P$ completely continuous?
(d) Classify the eigenvectors of $A$ according to eigenvalues of $P$.

\bigskip
\textbf{Solution.}

\subsection*{(a) $PA = AP$}

We need to show $(PAf)(t) = (APf)(t)$ for all $f \in L_2[0,1]$.

First, note that $a(t,t') = t_<(1-t_>) = \min(t,t')(1-\max(t,t'))$. Explicitly:
\[
a(t,t') = \begin{cases} t(1-t') & \text{if } t \le t', \\ t'(1-t) & \text{if } t' \le t. \end{cases}
\]

Compute $(PAf)(t) = (Af)(1-t)$:
\[
(Af)(1-t) = \int_0^1 a(1-t, t')\,f(t')\,dt'.
\]

Now compute $a(1-t, t')$. If $1-t \le t'$, i.e., $t' \ge 1-t$, then $a(1-t,t') = (1-t)(1-t')$. If $1-t > t'$, i.e., $t' < 1-t$, then $a(1-t,t') = t'(1-(1-t)) = t'\cdot t$.

So:
\[
a(1-t, t') = \begin{cases} t\,t' & \text{if } t' < 1-t, \\ (1-t)(1-t') & \text{if } t' \ge 1-t. \end{cases}
\]

Now compute $(APf)(t) = \int_0^1 a(t,t')\,f(1-t')\,dt'$. Substitute $u = 1-t'$, $du = -dt'$:
\[
(APf)(t) = \int_0^1 a(t, 1-u)\,f(u)\,du.
\]

We need $a(t, 1-u)$. If $t \le 1-u$ (i.e., $u \le 1-t$), then $a(t,1-u) = t(1-(1-u)) = tu$. If $t > 1-u$ (i.e., $u > 1-t$), then $a(t,1-u) = (1-u)(1-t)$.

So:
\[
a(t, 1-u) = \begin{cases} t\,u & \text{if } u \le 1-t, \\ (1-t)(1-u) & \text{if } u > 1-t. \end{cases}
\]

Comparing with $a(1-t, t')$ above (with $u$ in place of $t'$), we see $a(1-t, u) = a(t, 1-u)$ for all $t, u$. Therefore $(PAf)(t) = (APf)(t)$, so $PA = AP$. \hfill $\square$

\subsection*{(b) Self-adjointness, boundedness, and norm of $P$}

\textbf{Self-adjoint:}
\[
(Pf, g) = \int_0^1 f(1-t)\,\overline{g(t)}\,dt.
\]
Substituting $u = 1-t$:
\[
(Pf, g) = \int_0^1 f(u)\,\overline{g(1-u)}\,du = (f, Pg).
\]
So $P = P^\dagger$, i.e., $P$ is self-adjoint.

\textbf{Bounded:}
\[
\|Pf\|^2 = \int_0^1 |f(1-t)|^2\,dt = \int_0^1 |f(u)|^2\,du = \|f\|^2.
\]
So $\|Pf\| = \|f\|$, meaning $P$ is an isometry. Therefore $P$ is bounded with $\boxed{\|P\| = 1}$.

\subsection*{(c) Eigenvalues of $P$; complete continuity}

Since $P^2 f(t) = P(f(1-t)) = f(1-(1-t)) = f(t)$, we have $P^2 = I$. If $Pf = \lambda f$, then $P^2 f = \lambda^2 f = f$, so $\lambda^2 = 1$, giving $\lambda = \pm 1$.

Both eigenvalues are achieved:
\begin{itemize}
\item $\lambda = +1$: $f(1-t) = f(t)$ (symmetric functions), e.g., $f(t) = 1$, $f(t) = t(1-t)$.
\item $\lambda = -1$: $f(1-t) = -f(t)$ (antisymmetric functions), e.g., $f(t) = t - \tfrac{1}{2}$, $f(t) = \sin(2\pi t)$.
\end{itemize}

\textbf{$P$ is NOT completely continuous.} A completely continuous (compact) operator maps bounded sets to precompact sets. But $P$ is an isometry ($\|Pf\| = \|f\|$), and the identity operator on an infinite-dimensional space is not compact. Since $P^2 = I$, if $P$ were compact, then $I = P \circ P$ would be compact (composition of compact with bounded), which is impossible in infinite dimensions.

\subsection*{(d) Classification of eigenvectors of $A$}

Since $PA = AP$ and the eigenvalues of $P$ are $\pm 1$, by the simultaneous diagonalization theorem (Theorem~4.22), the eigenvectors of $A$ can be chosen to be simultaneous eigenvectors of $P$. That is, each eigenvector of $A$ can be chosen to be either symmetric ($Pf = f$, i.e., $f(1-t) = f(t)$) or antisymmetric ($Pf = -f$, i.e., $f(1-t) = -f(t)$) about $t = \tfrac{1}{2}$.

The known eigenvectors of $A$ (whose kernel is $t_<(1-t_>)$) are $x_n(t) = \sqrt{2}\sin(n\pi t)$ with eigenvalues $1/(n\pi)^2$. Indeed:
\begin{itemize}
\item For $n$ odd: $\sin(n\pi(1-t)) = \sin(n\pi)\cos(n\pi t) - \cos(n\pi)\sin(n\pi t) = \sin(n\pi t)$. So these are symmetric ($P$-eigenvalue $+1$).
\item For $n$ even: $\sin(n\pi(1-t)) = -\sin(n\pi t)$. So these are antisymmetric ($P$-eigenvalue $-1$).
\end{itemize}

\hfill $\blacksquare$

\end{document}
